\section{Discussion}

\label{sec:discussion}

\subsection{Decentralization Trade-offs}

The decentralized architecture introduces trade-offs compared to centralized repositories.
Query latency is higher (median 2.3 seconds vs.\ 0.4 seconds for GBIF) due to distributed index traversal.
Storage costs are approximately 3x higher than centralized cloud storage due to replication requirements.
Governance complexity increases for sovereign datasets, as community decision-making is slower than administrative access grants.

We argue these costs are justified by the gains in data integrity (immutable provenance), censorship resistance (no single point of control), and sovereignty (communities retain authentic authority).

\subsection{Adoption Barriers}

Three barriers limit broader adoption.
First, the token-based incentive system requires researchers to interact with cryptocurrency, creating friction for those unfamiliar with Web3 tools.
We are developing fiat on-ramps and institutional custody solutions to lower this barrier.
Second, existing data infrastructure investments create lock-in effects; institutions with established GBIF publishing workflows have limited incentive to adopt a parallel system.
Our GBIF synchronization adapter mitigates this by enabling dual publishing.
Third, the sovereign governance tier requires community capacity building that extends beyond technical deployment.

\subsection{Sustainability}

Long-term sustainability depends on the curation market generating sufficient incentives to maintain data quality and storage persistence.
Our economic model projects break-even at 50,000 active ZDOs, assuming current token valuation and storage costs.
Institutional partnerships for cold-tier storage provide a non-token-dependent persistence backstop.

% ============================================================
% 10. Conclusion
% ============================================================
